\subsection{平方secant正Pfaffianの主chartを全て閉じる}
\begin{theorem}\label{thm:h3-order2-22-secant-square-main-empty}
平方正top・対称secant負top・正Pfaffianの全 $z_5\ne0$ 領域にH3は存在しない。
$z_5=0$ 境界、非退化正topの正Pfaffian、及びrank2 secant境界はこの定理に含めない。
\end{theorem}
\begin{proof}
前節のshear被覆により
\[
 (z_0,z_1,z_2,z_3,z_4,z_5)=((1-t^2)/q,-t,q,0,t,q),\qquad q\ne0
\]
と書ける。
元の全四次関係から作った商汎関数のうち、六つだけを選ぶ。
この六関数の極は $q t(t^2-1)(q^2-4t^2+2)=0$ にしか存在せず、
正常式に対する指定 $10\times10$ minorは厳密に
\[
 -\frac{(2t-1)^3(2t+1)^3}
 {16q t^3(t^2-1)(q^2-4t^2+2)^2}
\]
である。従って表示した分母と $4t^2-1$ が非零なら正常核は零である。
全35関数を同時に扱うと現れる、これら六関数に不要な高次極を例外集合へ追加してはいない。

$t=1/2$ は前節で全正常pencilの因数分解により排除済み。
残る全例外を、次の五つの有理一パラメータ族で覆う：
\[
\begin{array}{c|c|c}
 q&t&\text{残る有限例外を含む式}\\\hline
 w&0&w^2+1\\
 w&1&1\\
 w&-1&1\\
 w&-1/2&w^2+1\\
 (w-2/w)/2&(w+2/w)/4&(w^2+2)(w^2+2w+2)P(w)
\end{array}
\]
最初の四行では $w\ne0$、最後の行では $w(w^2-2)\ne0$。
最後のcurveは $q^2-4t^2+2=0$ の全非零 $q$ 部分を、$w=2t+q$ により漏れなく覆う。
各行で元の関係から改めて四次汎関数を構成し、二つの正常10minorを計算する。
そのgcdと、実際に使う関数の極を合わせ、定義域外と既知の $t=1/2$ を除くと、
全残余は表の零点に含まれる。
これらの因子が次に述べる既知領域又は有限例外へ写ることも、剰余式で検証する。

最初と第四行の $q=\pm i$ はGaussian有理数体上で元の関係を再消去する。
$t=0$ では正常核は零。
$t=-1/2$ では正常核が一次元で、$q=i$ の全非零元は
\[
 (x+2y)\bigl(-x^2/2+xy+yx+2y^2-i(uv+vu)\bigr)
\]
に比例する。$q=-i$ はそのGalois共役であり、いずれも非domain切断となる。

最後の行の $w^2+2=0$ は、既に全値判定した $t=0,q^2=-2$ に写る。
$w^2+2w+2=0$ は上の $t=-1/2,q=\pm i$ に写る。
残る式は
\[
 P(w)=w^{12}-12w^{10}+28w^8+96w^6+112w^4-192w^2+64.
\]
これは $\mathbb Q$ 上既約である。
$\mathbb Q[\alpha]/(P(\alpha))$ 上で元の全四次関係を厳密消去すると、商次元35、
正自由三次32語への正常汎関数行列のrankは10となる。
従って正常三次核は零である。
計算は12次数体の元の多項式表現で行い、近似根によるrank判定ではない。
既約多項式から得た体の全埋込みが非零minorを保つので、これで十二根全てを覆う。

一般点・各curve・各有限点で用いる汎関数は、元の全四次関係を消す恒等式を持つ。
また、それらが作るrank10正常写像は正三次関係を消す。
従ってAS4候補の正三次商次元10と合わせれば正常核零となり、
特殊値でもgeneric商の基底を無条件に流用してはいない。
唯一の残存一次元核では上の実際の一次因子により商domain性を排除した。
これで $q\ne0$ の全主chartを閉じた。

検証器
\[
\texttt{verify\_cubic\_order2\_22\_secant\_square\_main.py}
\]
は選択六関数の極、二変数多項式行列式、五curveのgcd cover、全例外因子の写像、
二Gaussian対と唯一の因子、及び12次数体上の元関係消去を入力から再生する。
この主chartの排除を、残る $z_5=0$ や他の正topの完成と呼んではいない。
\end{proof}
